\section{Discussion}
\label{sec:discussion}

The document matters and its provenance does not: any of the three
documents takes the security checks passed by compiled artifacts from a
quarter or a third to well over four fifths, in both delivery modes, and
the one written without any repository does as well as the one anchored
to code. Functionality does not pay for it. Enforcement during the
attempts does not add to it: the judge leaves the rate where it is at
higher cost. What compiled document artifacts still fail depends on the
document's wording; the study repository records three revisions of the
agent's instructions, each removing failures the previous wording left
open, and is the place for that history.

\paragraph{Threats to validity.}
One task, one pair of games, one model as agent, generator and judge, ten
artifacts per arm. Each arm shares one acquired document, so the rates
measure generation variability given that document, not acquisition
variability. The five input-policy checks weigh heavily in the ten. The
agent's instructions were revised after examining failures on this task,
so they show that the rule can be written, not that it transfers.

\section{Conclusions}
\label{sec:conclusions}

Agent-written, repository-specific security context, given before the
model writes, raises the share of security checks a compiled artifact
passes from 22 to 33\% to 82 to 90\% without a functional cost, under one
response and under agentic delivery alike; a judge during the attempts
does not add to it. Whether the instructions transfer to another feature
and repository is the next question.
